% =========================================================================
% KAPITEL 5 -- SIEBEN ZUSAETZLICHE QUELLTEXTE UND EINE ZUSAMMENGELEGTE TABELLE
% 25.08.2026. Ergaenzt die acht bestehenden Ausschnitte; keiner faellt weg.
% Damit 15 Quelltexte: 3 in 5.1, 3 in 5.2, 3 in 5.3, 6 in 5.4.
%
% Alle Ausschnitte sind woertlich aus dem Repo, Stand 25.08.2026, nach den
% an diesem Tag nachgezogenen Kommentaren. Zeilennummern entsprechend.
%
% PLATZIERUNG
%   5.1 Absatz 1 -> lst:konstanten
%   5.1 Absatz 2 -> lst:gesperrt, lst:test-leakage
%   5.2 Absatz 3 -> lst:spatialjoin
%   5.2 Absatz 5 -> lst:joinhygiene
%   5.3 Absatz 2 -> lst:raster
%   5.4 Absatz 2 -> lst:test-folds  (steht schon, wandert nur von Absatz 16
%                                    hierher, wo die Aussage belegt wird)
%   5.4 Absatz 6 -> lst:streuung
%   tab:baselines ersetzt tab:baselines-menge und tab:baselines-struktur
% =========================================================================


% ---------------------------- AB HIER KOPIEREN ---------------------------

% ---------- 5.1 Absatz 1 -------------------------------------------------

\begin{lstlisting}[style=pythonstil,
  label={lst:konstanten},
  caption={Zeitraum und Ausschlüsse als benannte Konstanten
           (\texttt{prep/config.py}, Z.~83--97)}]
#   START  frueheste Periode mit vollstaendigen ACS-Merkmalen (#5, #11)
#   ENDE   letztes vollstaendiges Kalenderjahr (#12)
START = 201501
ENDE  = 202512

# Lag-Vorlauf (#23): ab START-VORLAUF aggregieren, Lags bilden, auf START
# zuschneiden. Ohne ihn beginnt die Regression erst 2016-01. 4.620 statt 4.200.
VORLAUF_MONATE = 12

# Parks ohne Wohnbevoelkerung (#19): 45 bis 507 Einwohner gegen 14.435 im
# Median. Jede Pro-Kopf-Groesse wird dort beliebig gross.
PARKGEBIETE = ["Golden Gate Park", "Lincoln Park", "Mclaren Park"]

# Warnschwelle fuer unvollstaendige Monate. Kein Filter - massgeblich ist ENDE.
VOLLSTAENDIGKEITS_SCHWELLE = 0.5
\end{lstlisting}


% ---------- 5.1 Absatz 2 -------------------------------------------------

\begin{lstlisting}[style=pythonstil,
  label={lst:gesperrt},
  caption={Gesperrte Ergebnisvariablen
           (\texttt{prep/config.py}, Z.~168--174)}]
# Ergebnisvariablen - NIEMALS Merkmal. Stehen erst nach dem Einsatz fest (#20).
ERGEBNISVARIABLEN = [
    "schaetzung_sachschaden_usd", "loeschfahrzeuge", "loeschkraefte",
    "rettungsdienst_einheiten", "alarmstufe", "antwortzeit_min",
    "zivile_tote", "zivile_verletzte", "flammenausbreitung_eingedaemmt",
    "ankunft_zeitpunkt",
]
\end{lstlisting}


\begin{lstlisting}[style=pythonstil,
  label={lst:test-leakage},
  caption={Prüfung auf Ergebnisvariablen in der erzeugten Datei
           (\texttt{tests/test\_aufbereitung.py}, Z.~277--285)}]
def test_keine_ergebnisvariablen():
    """Wichtigster Test des Strukturteils.

    Sachschaden, Loeschfahrzeuge, Alarmstufe und Antwortzeit stehen erst nach
    dem Einsatz fest. Rutscht eine dieser Spalten in den Merkmalssatz, sieht
    das Ergebnis gut aus und ist wertlos.
    """
    gefunden = [c for c in ERGEBNISVARIABLEN if c in klassifikation().columns]
    assert not gefunden, f"Ergebnisvariable(n) im Datensatz: {gefunden}"
\end{lstlisting}


% ---------- 5.2 Absatz 3 -------------------------------------------------

\begin{lstlisting}[style=pythonstil,
  label={lst:spatialjoin},
  caption={Räumlicher Verbund der Altquelle mit Abbruchschwelle
           (\texttt{prep/s1\_daten.py}, Z.~391--399)}]
    print("  Spatial Join (Deliktkoordinate -> Neighborhood-Polygon)...")
    punkte = gpd.GeoDataFrame(hist[["date"]].copy(),
                              geometry=gpd.points_from_xy(hist["x"], hist["y"]),
                              crs="EPSG:4326")
    joined = gpd.sjoin(punkte, neighborhoods_gdf(), how="left",
                       predicate="within")
    quote = joined["neighborhood"].notna().mean()
    assert quote >= 0.90, f"Match-Rate nur {quote:.1%} - Geometrie pruefen."
\end{lstlisting}


% ---------- 5.2 Absatz 5 -------------------------------------------------

\begin{lstlisting}[style=pythonstil,
  label={lst:joinhygiene},
  caption={Zeilenzahl vor und nach dem Verbund
           (\texttt{prep/s1\_daten.py}, Z.~643--647)}]
    vorher = len(base)
    base = base.merge(crime, left_on=["neighborhood", "year", "month"],
                      right_on=["neighborhood", "jahr", "monat"],
                      how="left").drop(columns=["jahr", "monat"])
    assert len(base) == vorher, "Crime-Join hat Zeilen dupliziert (1:n-Beziehung!)."
\end{lstlisting}


% ---------- 5.3 Absatz 2 -------------------------------------------------

\begin{lstlisting}[style=pythonstil,
  label={lst:raster},
  caption={Vollständiges Raster und Vorwärtsfüllen
           (\texttt{prep/s2\_datensaetze.py}, Z.~220--235)}]
    # Vollstaendiges Raster: Monate ohne Einsatz sind echte Nullen.
    idx = pd.MultiIndex.from_product(
        [agg["stadtteil"].unique(),
         range(int(agg["jahr"].min()), int(agg["jahr"].max()) + 1),
         range(1, 13)],
        names=["stadtteil", "jahr", "monat"])
    raster = agg.set_index(["stadtteil", "jahr", "monat"]).reindex(idx).reset_index()
    raster[ZIELGROESSE] = raster[ZIELGROESSE].fillna(0).astype(int)

    # Nur VORWAERTS fuellen. KEIN bfill: Rueckwaertsfuellen wuerde fehlende Werte
    # (z. B. akademikerquote vor ACS 2014) still mit ZUKUNFTSWERTEN imputieren -
    # Leakage (Decision Log #10). Echte NaN bleiben sichtbar.
    raster[_UEBERNOMMEN] = (raster.groupby("stadtteil")[_UEBERNOMMEN]
                                  .transform(lambda s: s.ffill()))
\end{lstlisting}


% ---------- 5.4 Absatz 6 -------------------------------------------------

\begin{lstlisting}[style=pythonstil,
  label={lst:streuung},
  caption={Zweistufige Mittelung über die zehn Wiederholungen
           (\texttt{vorpruefung/v1\_baselines.py}, Z.~181--206, gekürzt)}]
def _zweistufig(df: pd.DataFrame, schluessel: list[str],
                masse: list[str]) -> pd.DataFrame:
    g = df.groupby(schluessel, sort=False)
    z = g[masse].mean().add_suffix("_mean")
    z = z.join(g[masse].std().add_suffix("_std_folds"))

    # Streuung der 10 Wiederholungsmittel, nicht der 50 Laeufe: dieselben 29
    # Stadtteile in zehn Gruppierungen waeren zu optimistisch (R-5).
    je_wdh = df.groupby(schluessel + ["wiederholung"], sort=False)[masse].mean()
    z = z.join(je_wdh.groupby(schluessel, sort=False).std()
                     .add_suffix("_std_wiederholungen"))
\end{lstlisting}


% ---------- Ersetzt die beiden bisherigen Baseline-Tabellen --------------

\begin{table}[htbp]
    \centering
    \caption{Referenzwerte beider Stufen, Mittel über 50 Läufe; Streuung über die zehn Wiederholungsmittel}
    \label{tab:baselines}
    \footnotesize
    \singlespacing
    \setlength{\tabcolsep}{5pt}
    \begin{tabularx}{\textwidth}{@{}>{\raggedright\arraybackslash}p{3.9cm} c >{\raggedright\arraybackslash}X l l@{}}
        \toprule
        \textbf{Zielgröße} & \textbf{Stufe} & \textbf{Referenz} & \textbf{R\textsuperscript{2}} & \textbf{RMSE} \\
        \midrule
        Einsatzhäufigkeit & 1 & Gesamtmittelwert & $-0{,}744 \pm 0{,}325$ & $69{,}93 \pm 1{,}92$ \\
        Einsatzhäufigkeit & 2 & Poisson-Modell mit Offset & $0{,}542 \pm 0{,}082$ & $33{,}98 \pm 3{,}11$ \\
        Einsätze je 1.000 Einwohner & 1 & Gesamtmittelwert & $-1{,}054 \pm 0{,}875$ & $7{,}54 \pm 0{,}13$ \\
        Einsätze je 1.000 Einwohner & 2 & Poisson-Modell mit Offset & $0{,}367 \pm 0{,}261$ & $4{,}08 \pm 0{,}62$ \\
        \midrule
        & & & \textbf{Macro-F1} & \textbf{Trefferquote} \\
        \cmidrule(l){4-5}
        Überwiegende Einsatzart & 1 & Häufigste Klasse & 0,223 & 0,806 \\
        Überwiegende Einsatzart & 2 & Multinomiales Logit & $0{,}297 \pm 0{,}014$ & 0,584 \\
        \bottomrule
    \end{tabularx}
\end{table}

% ---------------------------- BIS HIER KOPIEREN --------------------------
